\subsection{کالبدشکافی ماژول \lr{samba\_m3}}

\subsubsection*{۱. فلسفه طراحی، پارادایم معماری و تمثیل ملموس}
معماری مقتدر \lr{samba\_m3} پیشگام ترکیب سه‌گانه مدل‌های حالت-فضا (\lr{Mamba SSM})، توجه موضعی با پنجره لغزان، و شبکه‌های پیش‌خور چندلایه است. این ترکیب هم‌افزا، حافظه استقرایی نامحدود \lr{SSM} را با وضوح فوق‌العاده توجه موضعی در هم می‌آمیزد.

\textbf{تمثیل ملموس و مثال کاربردی:}
سیستم محرکه خودروهای سه‌گانه هایبرید مدرن: موتور الکتریکی شتاب‌گیری‌های سریع و آنی در ترافیک شهری را تامین می‌کند (توجه موضعی با پنجره لغزان)، موتور احتراق داخلی کروز آرام و خطی در بزرگراه‌های طولانی بین‌شهری را هدایت می‌نماید (اسکن خطی سراسری مامبا)، و باتری تقویتی توان مضاعف برای سبقت‌گیری را مهیا می‌سازد (\lr{MLP}).

\subsubsection*{۲. فرمول‌بندی و استنتاج عمیق ریاضی \lr{(Mathematical Derivation)}}
توالی ورودی از لایه‌های متناوب فضای حالت بازگشتی و پنجره توجه محلی عبور می‌کند:
\begin{equation}
\mathbf{H}^{(l)} = \operatorname{MambaSSM}(\mathbf{X}^{(l)}), \quad \mathbf{Y}^{(l)} = \operatorname{SlidingWindowAttn}(\mathbf{H}^{(l)}, W)
\end{equation}
معادلات بازگشتی مامبا با اسکن موازی پیشوندی در زمان خطی حل می‌شوند:
\begin{equation}
\mathbf{h}_t = \bar{\mathbf{A}}_t \mathbf{h}_{t-1} + \bar{\mathbf{B}}_t \mathbf{x}_t, \quad \mathbf{y}_t = \mathbf{C}_t \mathbf{h}_t
\end{equation}
سپس پنجره محلی جزئیات ریز نحوی را بازآفرینی می‌کند تا پیچیدگی کل در سطح خطی $\mathcal{O}(N W + N d)$ بماند.

\subsubsection*{۳. کالبدشکافی الگوریتمی و جریان داده \lr{(Algorithmic Tensor Trace)}}
در پیاده‌سازی سورس پایتون در فایل ماژول \lr{samba\_m3.py}، پردازش در متد \lr{forward} به صورت مرحله‌به‌مرحله و گام‌به‌گام صورت می‌پذیرد:
\begin{enumerate}
\item دریافت تانسور ورودی و فرافکنی به فضای حالت پیوسته لایه مامبا.
\item اجرای اسکن خطی پیشوندی با عملگر شرکت‌پذیر بر روی کش سریع \lr{SRAM}.
\item اعمال لایه توجه موضعی با پنجره لغزان $W=128$ بر روی برآیند حالت‌ها.
\item ادغام با مسیر میان‌بر افزونه و عبور از لایه پیش‌خور \lr{SwiGLU}.
\item انتقال تانسور خروجی بازسازی‌شده به بلوک بعدی معماری هایبرید.
\end{enumerate}

\subsubsection*{۴. تحلیل همجوشی سخت‌افزاری و ترافیک حافظه \lr{(Hardware \& Roofline Analysis)}}
از دیدگاه مدل سقف کارایی (\lr{Roofline Model}) و معماری پردازنده‌های گرافیکی مدرن (مانند \lr{NVIDIA Ada / RTX 4070}}):
\begin{itemize}
\item \textbf{شکستن کامل سد حافظه در متون بلند:} به دلیل رفتار خطی مامبا، بافر حافظه با افزایش طول دنباله ثابت می‌ماند.
\item \textbf{بیشینه‌سازی توان عملیاتی هسته‌ها:} ترکیب اسکن سریع با پنجره کوچک توجه، هسته‌های پردازشی را در سقف توان نگه می‌دارد.
\end{itemize}

\subsubsection*{۵. ارزیابی تجربی، بنچمارک و توجیه علمی نتایج}
بر اساس آزمون‌های ثبت‌شده در خط لوله ارزیابی پایان‌نامه:
\begin{itemize}
\item \textbf{وضعیت رسمی آزمون کیفی:} \lr{PASS}
\item \textbf{دقت ثبت‌شده روی بنچ‌مارک:} 88.53 درصد
\item \textbf{شاخص ارزیابی آماری (\lr{F1 Score}):} 85.9 درصد
\item \textbf{اوج مصرف حافظه گرافیکی (\lr{Peak VRAM}):} 750.0 مگابایت (در برابر ۴۲۰۰ مگابایت مدل مبنا)
\item \textbf{نرخ توان عملیاتی (\lr{Throughput}):} 35,400 توکن بر ثانیه
\end{itemize}

\textbf{توجیه علمی و علت‌شناسی جذاب نتایج:}
دستیابی به دقت کم‌نظیر 88.53 درصد و سرعت 35,400 توکن/ثانیه با حفظ ردپای VRAM در 750.0 مگابایت (کاهش 82.1 درصدی نسبت به ۴۲۰۰ مگابایت پایه)، اثبات می‌کند که معماری هایبرید سامبا نقاط ضعف مدل‌های فضای حالت را پوشش داده و بهترین تعادل میان سرعت و حافظه را فراهم می‌آورد.

\subsubsection*{۶. دیاگرام معماری شکیل در لاتک \lr{(TikZ Architecture Flowchart)}}
جهت تبیین دقیق جریان داده تانسوری و نحوه تغییر شکل ماتریس‌ها، نمودار معماری این ماژول با استفاده از بسته قدرتمند \lr{TikZ} در محیط لاتک ترسیم شده است:

\begin{center}
\begin{tikzpicture}[scale=0.82, every node/.style={scale=0.82}, 
    node distance=1.1cm,
    auto,
    block/.style={rectangle, draw=blue!70, fill=blue!12, rounded corners=4pt, text centered, minimum height=0.85cm, font=\small},
    ssm/.style={rectangle, draw=purple!70, fill=purple!15, rounded corners=4pt, text centered, minimum height=0.85cm, font=\small},
    win/.style={rectangle, draw=orange!80, fill=orange!18, rounded corners=4pt, text centered, minimum height=0.85cm, font=\small},
    outblock/.style={rectangle, draw=teal!80, fill=teal!15, rounded corners=4pt, text centered, minimum height=0.85cm, font=\small}
]
    \node [block] (input) {ورودی $\mathbf{X}$};
    \node [ssm, right=1.0cm of input] (mamba) {مامبا \lr{SSM}};
    \node [win, right=1.0cm of mamba] (window) {پنجره لغزان};
    \node [outblock, right=1.0cm of window] (out) {خروجی هایبرید};

    \draw [->, thick, draw=blue!80] (input) -- (mamba);
    \draw [->, thick, draw=purple!80] (mamba) -- (window);
    \draw [->, thick, draw=orange!80] (window) -- (out);
\end{tikzpicture}
\end{center}

\subsubsection*{۷. پیوست تصویری و پرامپت تولید هوش مصنوعی}
برای ترسیم تجسم سه‌بعدی و گرافیک صنعتی این مکانیزم با ابزارهای هوش مصنوعی، پرامپت انگلیسی زیر با دقت تدوین شده است:

\vspace{0.2cm}
\begin{latin}
\footnotesize
\textbf{AI Technical Visualization Prompt:}\\
"A 3D scientific visualization of the Samba-M3 tri-hybrid architecture. A continuous cyan ribbon representing Mamba SSM state-space scan intertwining with an orange geometric ring representing local sliding-window attention, high-tech industrial rendering, 8k resolution --ar 16:9"
\end{latin}

\begin{figure}[H]
    \centering
    \framebox[0.9\textwidth]{
    \begin{minipage}{0.85\textwidth}
    \vspace{1.8cm}
    \centering
    \textbf{[کادر الحاق تصویر گرافیکی سه‌بعدی ماژول \lr{samba\_m3}]}\\[0.2cm]
    \small (بر اساس پرامپت فنی فوق تولید و در این کادر جایگزین می‌شود)
    \vspace{1.8cm}
    \end{minipage}
    }
    \caption{شماتیک و جریان داده تانسوری در ماژول \lr{samba\_m3}}
    \label{fig:samba_m3}
\end{figure}

\vspace{0.5cm}
\hrule
\vspace{0.5cm}
